\documentclass[11pt,a4paper]{article}
\usepackage[margin=24mm]{geometry}
\usepackage{amsmath,amssymb,booktabs,microtype,graphicx}
\usepackage[hidelinks]{hyperref}
\newcommand{\dd}{\mathrm d}
\title{Route G2-X: An Externally Anchored Xenon Feasibility Test}
\author{A real response matrix; conditional event yield, not a built source}
\date{23 September 2026}
\begin{document}
\maketitle
\begin{abstract}
We conditionally select the uniform $X$-only Higgs portal and fold the
existing driven recoil population through the published XENON1T S2-only
nuclear-recoil response. A Higgs-threshold-motivated scale interval is
tested against a direct invisible-width constraint and the retained
action's vacuum condition. Unknown incident flux remains an explicit
axis of feasibility, not an invented experimental capability. A useful
negative result follows: throughout the retained recoil-energy window,
the scalar scattering kernel factorizes, so this energy-only channel
cannot reproduce the old two-line momentum classifier.
\end{abstract}

\section{Physical choices and provenance}
This is a conditional phenomenological branch, not a prediction of the
radius. Retain $\alpha=.25$, $g=.05$, the selected $m=0,|j|=1$ source
population, $\sigma_i=.1\Lambda$, and the activated drive
$\epsilon=.5$, $\Omega=.2\Lambda$. Choose
\begin{equation}
 1\leq\Lambda=R^{-1}\leq30\ {\rm GeV},\qquad
 \mu_j=\Lambda\sqrt{25+j^2},\qquad m_X=\mu_{\pm1}=\sqrt{26}\Lambda.
\end{equation}
The range spans the physical Higgs pair threshold
$\Lambda_h=m_h/10=12.508$ GeV and $X_0$ masses 5--150 GeV.
These boundaries are a coverage choice, not a measured dark mass or a
fit to known particles. Use $m_h=125.08$ GeV, $v=246$ GeV, and
$\Gamma_h^{\rm SM}=4.10$ MeV from the currently served PDG review
\cite{pdg}; $\Gamma_h^{\rm SM}$ is a theoretical SM prediction.
The coupling domain $10^{-5}\leq\lambda_X\leq1$ is a declared scan
range, not a confidence interval or a perturbativity theorem.

The target is natural liquid xenon, using the \emph{actual} official
XENON1T 2019 S2-only analysis/2020 response release \cite{xenon,release}.
We choose the older, directly usable response product instead of calling
a newer halo-WIMP exclusion a constraint on this boosted source.
Its source commit is
\texttt{5a364bc8709f2561e5a013ddea6993a5a7c8e313}.
Both CSV files and their hashes, original documentation and attribution
are retained locally. The response includes event-selection losses
averaged over the full target. The raw search exposure is
\begin{equation}
 {\cal E}=356770\ {\rm kg\,day}.
\end{equation}
No separate fiducial-efficiency factor is multiplied in. This reuse
assumes approximately uniform, thin-target illumination, the source
COM frame at rest with respect to the target, and the released
energy-only response appropriate to these recoils. It is not a
calibration of a new narrow-beam installation.

The released NR matrix extends to hypothetical response below the
analysis threshold. We explicitly retain true recoil energy
$0.7\leq E_R\leq50$ keV and only full S2 bins inside 150--3000 PE.
The actual lower retained edge is approximately 150.027 PE; the bin
straddling 150 PE is dropped without assuming a sub-bin shape.
The response is used as supplied, not replaced by a Gaussian. Energies
outside this interval are unmodeled, not asserted to produce zero
physical response. No response covariance or new background likelihood
is invented.

\section{Constraints from the same candidate interaction}
At a declared common tree matching prescription, select
\begin{equation}
 \Delta{\cal L}=-\lambda_X H^\dagger H\sum_j|x_j|^2,
 \quad \lambda_X\geq0,\quad \lambda_\Phi=0.
\end{equation}
This is a four-dimensional EFT option. The $X$ physical intercept is
held fixed by explicit matching,
\begin{equation}\label{eq:vacuum}
 M_{D,0}^2=25\Lambda^2-\lambda_Xv^2/2\geq0,
 \qquad\lambda_X\leq\lambda_{\rm vac}=50\Lambda^2/v^2 .
\end{equation}
With no $|X|^4$ term in the retained action, violating this condition
gives a runaway along $H=\Phi=0$, constant $X$. Its exclusion is a
theory failure, distinct from an experimental exclusion. The analogous
claim for a negative $\Phi$ bare mass needs its holonomy gradient gap;
we do not import the $X_0$ argument into that sector.

The published ATLAS combination gives $B_{\mathrm{inv}}<.107$ at 95\% CL,
assuming SM Higgs production \cite{atlas}. With unchanged visible
couplings and all produced dark modes invisible, the permitted new
width is
\begin{align}
 \Gamma_{\max}&=\frac{.107}{1-.107}\Gamma_h^{\rm SM},\\
 \Gamma_X&=\frac{\lambda_X^2v^2}{16\pi m_h}B_X(\Lambda),\qquad
 B_X=\sum_{j:2\mu_j<m_h}\sqrt{1-4\mu_j^2/m_h^2},\\
 \lambda_X&\leq\lambda_h(\Lambda)
 =\sqrt{\frac{16\pi m_h\Gamma_{\max}}{v^2B_X(\Lambda)}}
 \quad(B_X>0).
\end{align}
Every signed label is one complex field with a distinct
particle--antiparticle final state. Above $\Lambda_h$ this particular
on-shell decay supplies no bound, not a declaration that all collider
constraints disappear. The map labels this a conditional tree recast
of a verified 2023 input, not the latest exhaustive global exclusion.

At the matched static tree spectrum, spontaneous contact decay is
closed. For $X_j\to X_l+\Phi_n+\bar\Phi_m$ with $j=l+n-m$,
the triangle inequality gives
\begin{equation}
 \mu_l+m_n+m_m
 \geq\sqrt{(M_D+2M_\Phi)^2+j^2\Lambda^2}>\mu_j.
\end{equation}
The particle/antiparticle holonomies cancel. A uniform Higgs cannot
change $j$. Collider escape additionally requires negligible material
scattering; it is an explicit recast assumption. The collider is not
placed in the hypothetical pump: this constraint uses the drive-off
vacuum, while the source forecast retains the drive.

Mixed $g\lambda_X$ loops permit a $\Phi$ portal counterterm. Setting it
to zero at matching is not an all-orders protection. The cutoff must
control the Higgs energy and the modes retained in each application;
neither a KK species-count heuristic nor small four-dimensional
couplings proves that control. We keep this bounded tree feasibility
test conditional, rather than initiate another full matching program.

\section{Finite-speed nuclear recoil and the measured response}
Use natural xenon number fractions from CIAAW \cite{xe}, molar mass
$\bar M=131.293$ g/mol, and the explicit approximation
$m_A=A\times.93149410242$ GeV. Let $m_N=.939$ GeV and
$f_N=.308\pm.018$ \cite{fn}; its effective two-nucleon contribution is
not added a second time. We approximate the coherent nuclear scalar
current by its static normalization and the Helm form factor
\cite{helm},
\begin{equation}
 F_A(q)=\frac{3j_1(qr_A)}{qr_A}e^{-(qs_0)^2/2},\quad
 r_A^2=c_A^2+7\pi^2a^2/3-5s_0^2,
\end{equation}
with $c_A=1.23A^{1/3}-.60$ fm, $a=.52$ fm and $s_0=.9$ fm.
Arguments use $\hbar c=.1973269804$ GeV fm. This is not a complete
finite-$q$ nuclear many-body calculation. The admitted window has
$q\lesssim.113$ GeV; no coherent formula is extrapolated to the much
larger full recoil endpoints.

For incident momentum $p_X$ in the target rest frame,
\begin{align}
 s&=m_X^2+m_A^2+2m_A\sqrt{m_X^2+p_X^2},\qquad
 {\cal K}(s,m_X^2,m_A^2)=4m_A^2p_X^2,\\
 \overline{|{\cal M}_A|^2}
 &\simeq\frac{4\lambda_X^2 f_N^2m_N^2A^2m_A^2}
 {(m_h^2+2m_AE_R)^2}F_A^2(\sqrt{2m_AE_R}),\\
 \frac{\dd\sigma_A}{\dd E_R}
 &=\frac{\lambda_X^2f_N^2m_N^2A^2m_A}
 {8\pi p_X^2(m_h^2+2m_AE_R)^2}
 F_A^2(\sqrt{2m_AE_R})\,\mathbf1_{0<E_R<E_R^{\max}},\\
 E_R^{\max}&=\frac{2m_Ap_X^2}{s}.
\end{align}
The incident kinematics are not replaced by a halo-speed distribution.
In the NR, zero-form-factor limit, integrating gives
$\sigma_A=\lambda_X^2f_N^2m_N^2A^2\mu_{XA}^2/
(4\pi m_h^4m_X^2)$, which supplies a normalization check.

Let $r=(s_j,q,p_i)$ label the existing selected source mixture. Here
$p_i$ is the dimensionless G2 variable and $p_X=\Lambda\hat p_f(p_i)$;
the corresponding physical incoming width is $.1\Lambda$. Then
\begin{equation}
 \dd P_r=Z_{\rm pair}^{-1}w_{l_{s_j}}f_{.1}(p_i)
 K_{l_{s_j},0,q}(p_i)\dd p_i,\qquad l_+=0,\ l_-=-2.
\end{equation}
These are event-weighted incident \emph{flux fractions} at the sensor
only if transmission is common; unknown source luminosity and transport
are not thereby computed. With $R_b(E_R)$ the released probability
of an S2 bin $b$, including losses, define
\begin{equation}
 \bar\sigma_b=\sum_A\xi_A\int\dd P_r
 \int_{.7\,\mathrm{keV}}^{50\,\mathrm{keV}}
 R_b(E_R)\frac{\dd\sigma_A(p_X(r))}{\dd E_R}\dd E_R.
\end{equation}
Use the supplied energy-bin edges with piecewise-constant response and
integrate the smooth cross section within each clipped bin. No
independent efficiency correction or repeated velocity factor is used.
The unit conversion is $1\ {\rm GeV}^{-2}=3.893793721\,10^{-28}$ cm$^2$.

For stationary, spatially uniform incident selected-particle flux
$\varphi_X$ at the detector in cm$^{-2}$s$^{-1}$,
\begin{equation}\label{eq:flux}
 N_b=\varphi_X\,86400\,{\cal E}\frac{1000N_A}{\bar M}\bar\sigma_b,
 \qquad
 \varphi_3=\frac3{86400\,{\cal E}(1000N_A/\bar M)
                      \sum_{b\in\mathrm{ROI}}\bar\sigma_b}.
\end{equation}
Here $N_A$ is Avogadro's constant, not a nuclear mass number.
Three expected signal events is a transparent normalization target,
\emph{not} a discovery threshold, a 95\% sensitivity, or a fit to
observed events. The source flux is missing, so data cannot yet yield
a two-parameter xenon exclusion in $(\Lambda,\lambda_X)$.

\section{Why this channel loses the old sign information}
The exact kinematic lower bound over the entire source support is
$p_X\geq .8091444644\Lambda$, attained as $p_i\to0$ in the
$j=+1,q=-1$ branch. It follows from monotonicity of outgoing two-body
momentum with available energy. At fixed dimensionless $\hat p,\hat m$
and nuclear mass, write $\hat E=\sqrt{\hat m^2+\hat p^2}$ and
$D=m_A^2+2m_A\hat E\Lambda+\hat m^2\Lambda^2$. Then
\begin{equation}
 \frac{\dd E_R^{\max}}{\dd\Lambda}
 =\frac{4m_A\hat p^2\Lambda(m_A^2+m_A\hat E\Lambda)}{D^2}>0.
\end{equation}
Checking $\Lambda=1$ for every isotope therefore bounds the whole
continuous interval, not just sampled scales. The minimum endpoint
is 9.543 MeV, above 50 keV throughout the range.
Consequently the endpoint indicator is always one in the retained
window, and the cross section factorizes:
\begin{align}
 \frac{\dd\sigma_A}{\dd E_R}=
 \frac{\lambda_X^2}{p_X^2}H_A(E_R),&\qquad
 \bar\sigma_{s_jb}=\lambda_X^2 C_{s_j}D_b,\\
 C_{s_j}=\int_{s_j}\frac{\dd P_r}{p_X(r)^2},&\qquad
 D_b=\sum_A\xi_A\int_{.7\,\mathrm{keV}}^{50\,\mathrm{keV}}
 R_b(E_R)H_A(E_R)\dd E_R.
\end{align}
Any common response $R_b(E_R)$ preserves this factorization. Thus
\begin{equation}
 P(b|j=+,\mathrm{accepted})=P(b|j=-,\mathrm{accepted}),\qquad
 P(j|b,\mathrm{accepted})=\frac{C_j}{C_++C_-}.
\end{equation}
The optimal energy-only classifier is the majority guess, with error
$\min(C_+,C_-)/(C_++C_-)$. Better S2 resolution cannot restore the
missing momentum dependence within this approximation and window.
The original 2.20\% error was for a directly measured momentum
estimator, not a nuclear recoil deposit. The negative result is about
this specific readout, not about the existence of the mode conversion.
Response leakage from unmodeled energies, a direction-sensitive response,
or additional observations require a different calculation.

\section{Numerical feasibility, with the source left unknown}
Independent source integration gives, with $\Lambda$ in GeV,
\begin{equation}
 (C_+,C_-)=\frac{(.74075482,.12115857)}{\Lambda^2},\qquad
 P(j=\pm\mid\mathrm{accepted})=(.85943069,.14056931).
\end{equation}
Hence the best S2-only sign error is $14.0569\%$, exactly the minority
prior within this model; the conditional mutual information is zero.
The source quadrature retains its existing $z\leq10$ tail convention;
the kinematic endpoint proof does not rely on that numerical truncation.
The accepted-window normalization is
\begin{align}
 \bar\sigma_{\rm acc}&=7.99989\,10^{-39}\,
       \frac{\lambda_X^2}{(\Lambda/{\rm GeV})^2}\ {\rm cm}^2,\\
 N_{\rm acc}&=.00113108535\,
       \frac{\varphi_X}{1\ {\rm cm}^{-2}{\rm s}^{-1}}\,
       \frac{\lambda_X^2}{(\Lambda/{\rm GeV})^2}.
\end{align}
The unit-coupling, unit-scale point is a normalization convention,
not itself an allowed benchmark. For each scale use
$\lambda_{\max}=\min(1,\lambda_{\rm vac},\lambda_h)$, omitting
$\lambda_h$ when no Higgs channel is open. Table~\ref{tab:flux}
reports the smallest required flux within that conditional domain.

\begin{table}[ht]
\centering\small
\begin{tabular}{rrrrl}
\toprule
$\Lambda$ [GeV]&$m_X$ [GeV]&$\lambda_{\max}$&
$\varphi_3$ [cm$^{-2}$s$^{-1}$]&Active condition\\
\midrule
1&5.099&.000723&$5.07\,10^9$&Higgs width\\
5&25.495&.00176&$2.15\,10^{10}$&Higgs width\\
10&50.990&.00382&$1.82\,10^{10}$&Higgs width\\
12&61.188&.00857&$5.21\,10^9$&Higgs width\\
12.508&63.779&.129&$2.48\,10^7$&Bare-mass condition\\
30&152.971&.744&$4.32\,10^6$&Bare-mass condition\\
\bottomrule
\end{tabular}
\caption{Three expected recoils, not discovery sensitivity. The narrow
turn near $m_h/10$ is the closing on-shell decay threshold, not a
measured preferred radius or an optimized viable point.}
\label{tab:flux}
\end{table}

The quoted $f_N$ uncertainty alone multiplies yields by .8865--1.1203,
or required flux by .8926--1.1280. This is not a total error band:
nuclear, response, source and EFT uncertainties remain unquantified.
The new verifier passes 573/573 checks; all 3056 preceding Route-G
checks also pass unchanged (3629 total). These test implementation
and declared mathematical identities, not empirical validation of $X$.
The dense mesh in Fig.~\ref{fig:map} only renders these analytic
formulas; it is not another source or Gaussian-resolution scan.

\section{Decision and remaining physical inputs}
The accompanying figure separates the conditional ATLAS exclusion,
the retained-action vacuum failure, and the flux needed for three
accepted recoils. No white or unhatched region is called experimentally
viable: other constraints, production, transport and theory errors are
not closed. The response-derived result concerns a calibrated component
of an instrument, not the source, full sensor system, or sign classifier.

The energy choice also exposes the drive requirement:
$\Omega=.2$--6 GeV corresponds to ordinary frequency
$\nu=\Omega/h\simeq4.84\,10^{22}$--$1.45\,10^{24}$ Hz.
No physical apparatus that modulates the quartic this way has been
specified. It remains a prescribed reservoir, not an ordinary laboratory
oscillator; an actual source or a justified quantum energy-exchange
process is needed before hardware feasibility can be claimed.

Two productive next choices remain. If only an excess of accepted
recoils is required, specify production and at-detector flux, then
compare with \eqref{eq:flux} and an actual background likelihood.
If the purpose is to distinguish $j=\pm1$, change the observable:
independent timing/momentum information or a recoil endpoint/directional
measurement must be developed. Do not try to recover the missing
information by fitting a narrower Gaussian to this S2-only dataset.

\begin{thebibliography}{9}
\bibitem{pdg} PDG, \emph{Electroweak Model and Constraints on New Physics},
Sec.~10.4.2, currently served April 2026 revision,
\href{https://pdg.lbl.gov/2025/reviews/rpp2025-rev-standard-model.pdf}{review};
\href{https://pdg.lbl.gov/2025/reviews/rpp2025-rev-phys-constants.pdf}{physical constants}.
\bibitem{xenon} XENON Collaboration, \emph{Light Dark Matter Search with
Ionization Signals in XENON1T}, Phys. Rev. Lett. 123, 251801 (2019),
\href{https://arxiv.org/abs/1907.11485}{arXiv:1907.11485}.
\bibitem{release} XENON Collaboration, \emph{XENON1T S2-only data release}
(2020), \href{https://github.com/XENON1T/s2only_data_release}{official repository}.
\bibitem{atlas} ATLAS Collaboration, Phys. Lett. B 842, 137963 (2023),
\href{https://arxiv.org/abs/2301.10731}{arXiv:2301.10731}.
\bibitem{fn} M.~Hoferichter et al., Phys. Rev. Lett. 119, 181803 (2017),
\href{https://arxiv.org/abs/1708.02245}{arXiv:1708.02245}.
\bibitem{helm} J.~D.~Lewin and P.~F.~Smith, Astropart. Phys. 6, 87 (1996),
\href{https://hepwww.pp.rl.ac.uk/groups/ukdmc/pub/papers/journal/app6-87.pdf}{Eqs.~(4.7), (4.10--11)}.
\bibitem{xe} CIAAW, \href{https://ciaaw.org/xenon.htm}{Xenon isotope
abundances and standard atomic weight}.
\end{thebibliography}
\clearpage
\begin{figure}[p]
\centering
\includegraphics[width=\textwidth]{route_g/output/figures/g2_xenon_feasibility.pdf}
\caption{Conditional feasibility and the sign-readout obstruction.
Color encodes a required, not a supplied, flux. Panel B uses the
unchanged public response matrix and retains its bin structure.}
\label{fig:map}
\end{figure}
\end{document}
